\subsection{Приближение статистик}

\subsubsection{Математическое ожидание}

Рассмотрим реализованную выборку \( X \):

\begin{center}
\begin{tabular}{|c|c||c|c||c|c||c|c|}
\hline
$i$ & $x_i$ & $i$ & $x_i$ & $i$ & $x_i$ & $i$ & $x_i$ \\ \hline \hline
1 & 2.89 & 6  & 3.32 & 11 & 3.09 & 16 & 3.30 \\ \hline
2 & 3.14 & 7  & 3.01 & 12 & 2.66 & 17 & 3.07 \\ \hline
3 & 3.12 & 8  & 2.55 & 13 & 3.21 & 18 & 3.42 \\ \hline
4 & 3.34 & 9  & 3.25 & 14 & 2.95 & 19 & 3.05 \\ \hline
5 & 2.89 & 10 & 3.24 & 15 & 3.38 & 20 & 3.03 \\ \hline
\end{tabular}
\end{center}

Для оценки математического ожидания реализованной выборки воспользуемся выборочным моментом первого порядка \textit{(другими словами, средним выборочным)}:
\[
    \overline{X}=\frac{1}{n}\sum_{i=1}^n X_i
\]

Посчитаем для реализованной выборки при помощи скрипта на \textit{Python}:
\begin{lstlisting}
sample = [2.89, 3.14, 3.12, 3.34, 2.89, 3.32, 3.01, 2.55, 3.25, 3.24, 3.09, 2.66, 3.21, 2.95, 3.38, 3.30, 3.07, 3.42, 3.05, 3.03]
n, M = len(sample), 0
for i in range(n):
  M += sample[i]
M /= n
print("Математическое ожидание:", M)
\end{lstlisting}

Итого:
\[
    \overline{x}=3.0955\implies\boxed{
        MX\approx 3.0955
    }
\]

\subsubsection{Дисперсия}

Заметим, что объём выборки \( n=20<100 \) --- довольно малое значение. Значит, для оценки дисперсии необходимо воспользоваться \textit{несмещённой выборочной дисперсией}:
\[
    S_0^2=\dfrac{1}{n-1}\sum_{i=1}^n(X_i-\overline{X})^2,
\]

где \( \frac{1}{n-1} \) --- поправка Бесселя.

Посчитаем выборочные моменты второго порядка при помощи скрипта:
\begin{lstlisting}
D = 0
for i in range(len(sample)):
  x_i = sample[i]
  x2 = (x_i - M)**2
  D += x2
  print(f"x_{i + 1}".rjust(4), end=" | ")
  print(f"{x_i:.2f}", end=" | ")
  print(f"{x2:.5f}")
D /= n - 1
print("Дисперсия:", f"{D:.5f}")
\end{lstlisting}

Заполним таблицу с промежуточными значениями для наглядности:

\begin{table}[h!]
\centering
\begin{tabular}{|c|c|c||c|c|c|}
\hline
$i$ & $x_i$ & $(x_i - \bar{x})^2$ & $i$ & $x_i$ & $(x_i - \bar{x})^2$ \\ \hline \hline
1 & 2.89 & 0.04223 & 11 & 3.09 & 0.00003 \\ \hline
2 & 3.14 & 0.00198 & 12 & 2.66 & 0.18966 \\ \hline
3 & 3.12 & 0.00060 & 13 & 3.21 & 0.01311 \\ \hline
4 & 3.34 & 0.05978 & 14 & 2.95 & 0.02117 \\ \hline
5 & 2.89 & 0.04223 & 15 & 3.38 & 0.08094 \\ \hline
6 & 3.32 & 0.05040 & 16 & 3.30 & 0.04182 \\ \hline
7 & 3.01 & 0.00731 & 17 & 3.07 & 0.00065 \\ \hline
8 & 2.55 & 0.29757 & 18 & 3.42 & 0.10530 \\ \hline
9 & 3.25 & 0.02387 & 19 & 3.05 & 0.00207 \\ \hline
10 & 3.24 & 0.02088 & 20 & 3.03 & 0.00429 \\ \hline
\end{tabular}
\caption{Расчет дисперсии}
\end{table}

Итого:

\[ 
    S_0^2=0.05294
    \implies\boxed{
        DX\approx 0.05294
    }
\]